\chapter{When something looks wrong}\label{ch:trouble}
\begin{goals}
\textbf{In this chapter} you find the symptom you are seeing, why wuffle does it, and what to do. Every
case here was reproduced on the program itself; none is a guess.
\end{goals}

The cases are grouped as the book is: the compass, opening data, the nets, the console, and projects.
Each starts with what you see.

\section{The compass}\label{sec:trouble-field}

\subsection{The compass reads nothing}

The status line says \ui{no orientation yet — tilt the phone (or type a measurement)} when you press
\ui{◎}. The browser is giving wuffle no orientation:

\begin{itemize}
\item on an iPhone or iPad, Safari never does, because wuffle doesn't ask its permission
(\cref{sec:compass-browsers});
\item a page served from a plain \file{http://} address gets no sensors; open the site, which is https;
\item a computer has no orientation sensor at all.
\end{itemize}

Typed measurements work everywhere (\cref{sec:manual-entry}).

\subsection{The status line says the compass is low or unreliable}

Inside the lead-acid shell, wuffle shows the magnetometer's own calibration report:
\ui{compass low — wave the phone in a figure-8 to calibrate}. A phone's compass drifts near steel,
vehicles, power lines and in pockets with keys. Wave it through a figure of eight a few times, away from
anything magnetic, until the status says \ui{compass calibrated}.

\subsection{The readings are magnetic, not true}

The log's \col{north} column says \code{magnetic} and \col{declination\_deg} is empty. wuffle corrects to
true north only from a position (\cref{sec:true-north}), and it gets a position only inside the shell,
after the first GNSS fix. Readings taken before the fix, or outside the shell, stay magnetic, and say so
on every row. Correct them before mixing them with true-north data: add the declination for the place and
date to each azimuth.

A typed measurement says \code{as-typed}: wuffle takes it as you wrote it.

\subsection{The log was empty when I came back}

The compass keeps its log on the phone, in the browser's own storage, after every measurement, and
brings it back when the page opens again: the message line says \ui{kept from last time: 14 measurements
not yet cleared — publish to send them}. A closed tab, a reload, or the phone closing the browser to
save memory loses nothing.

What does lose it:
\begin{itemize}
\item \ui{clear}, tapped twice. The first tap only asks (\ui{tap clear again to delete all 14
measurements — publish first to keep them}); the second, within three seconds, empties the log for
good.
\item clearing the browser's data for the page, or a private window, which keeps nothing once it closes.
\item a different address: the log is kept per site, so the same compass opened from a file and from
gentropic.org/wuffle has two logs.
\end{itemize}

Publish as you go all the same: \ui{publish →} at each station, or every few, costs nothing, and a later
publish writes the whole log again. A published log can be taken back into the compass inside the shell
(share the CSV to wuffle), or opened on the desk.

\subsection{A line I typed went in as a plane}

\code{200/20} is a plane, \code{200→20} and \code{200>20} are lines. With the \ui{line} mode on,
\code{200/20} is taken as a line too. Delete the row with its \ui{×} and type it again.

\section{Opening data}\label{sec:trouble-data}

\subsection{The file opened, but with the wrong columns}

Look at the table's heading. When wuffle had to guess which columns hold the attitude, it says so, and
how it read them:

\shot{it-guess}{A file whose columns, \col{azimuth} and \col{inclination}, aren't names wuffle knows: it
read the first two numeric columns as dip direction and dip, and says so in the heading.}

The guess is the first two numeric columns, as dip direction then dip, swapped when the second runs past
90. If that is wrong for your file, rename its header to names wuffle knows and open it again:
\col{dipdir} and \col{dip} for planes, \col{strike} and \col{dip} for the right-hand rule, \col{trend}
and \col{plunge} for lines (\cref{sec:data-roles} lists every name it recognises). A pair that mixes
kinds, such as \col{azimuth} (a trend) with \col{inclination} (a dip), is no pair, and falls to the
guess.

\subsection{Some rows aren't drawn}

A row whose attitude can't be read (a dip of \code{abc}, a dip over 90, an empty cell) stays in the
table, with \code{—} in the first column, and is left off every diagram. The layer's properties count
them: \ui{tocher · planes · 198 of 200 rows read}. Fix the cells in the source file, or in the table
(\cref{sec:table}), and they are drawn.

\subsection{My file of \code{045→30} turned into two datasets}

A single column holding attitudes is split by the arrow: rows with one are lines, the rest planes, and a
column with both opens as \code{<name> · planes} and \code{<name> · lines}. That is deliberate: planes
and lines don't share a statistics card.

\subsection{The log's photos didn't come in}

Open the photos \emph{with} the log, in the same pick or the same drop. wuffle takes the photos the
log's \col{photo} column names, and says how many it found: \ui{(2 named in the log weren't in the
drop)}. Photos opened on their own are ignored, with the same advice in the status line.

\section{The nets}\label{sec:trouble-nets}

\subsection{A note in the net's corner}

What the net couldn't draw, it says, in yellow at the bottom left:

\shot{it-note}{Two layers of one dataset, both with contours on: the net draws one set, and says
whose.}

\begin{description}
\item[\ui{contours: one set per net — “…” isn’t contoured; the contours are “…”’s}] A net draws the
contours of the first contoured layer only. Contour the other dataset on a net of its own
(\menu{Diagram>New net}). When both layers show the same dataset, the note names it twice; a layer
drawn from the console with a link name of its own
(\code{ln -s ../../../data/tocher.csv "diagrams/Net 1/layers/second look.csv"}) carries that name, and
the note uses it.
\item[\ui{right dihedra aren’t drawn on a rotated net (Reset rotation to see them)}] The right
dihedra are computed on the unrotated net. Right-click the net's name and choose \ui{Reset rotation}
(\cref{sec:rotate-reset}).
\end{description}

\subsection{Contours are switched on but nothing is drawn}

Contours need at least three readable measurements. A layer with fewer draws none, and the net's corner
says so: \ui{contours: “east” has 2 measurements — contours need three}. Check the row count beside the
layer in the tree, and the rows it couldn't read (\cref{sec:data-read}).

\subsection{The mean sits where there is no data}

On a girdle (the Tocher sample is one) the Fisher mean of the poles lands wherever the spread happens to
balance, which need not be where the poles gather, and its cone is wide: on Tocher, $\kappa$ is 3.1 and
$\alpha_{95}$ 6.8°. That is the mean of a distribution a mean doesn't describe. Turn on the eigenvectors: V3, the girdle's axis, and the V3 great
circle, the girdle itself, describe it (\cref{sec:eigen}).

\subsection{A cone selected nothing, or the wrong poles}

\code{select cone 040/50 20} is a cone about the \emph{pole} of the plane 040/50, which is
\code{220→40}: on the Tocher sample both select 18 poles. \code{040→50} is a different direction
altogether (6 poles). Type a plane when you mean \enquote{planes like this one}, a line when you mean
\enquote{poles near this direction}. The cone tool on the net, dragged from the centre you want, has no
such ambiguity (\cref{sec:shapes}).

\subsection{A rose has petals at both ends}

A strike rose always does: a strike of 030 is the same line as 210. A dip-direction or trend rose has
them only when \ui{Count both ends (axial)} is ticked in its menu (\cref{sec:rose}); untick it.

\subsection{Labels cover the net}

Labels name every measurement of a layer, so 200 of them are a blur. Select the few to name, make them a
dataset (\cref{sec:select-dataset}), draw that, and label only it (\cref{sec:labels}).

\section{The console}\label{sec:trouble-console}

\shot{it-redirect}{The redirect trap: in \file{/home}, \code{select cone 150>60 10} made a file called
\file{60}; quoted, the same attitude reaches the command.}

\subsection{A command with \code{>} failed, or made a file}

In the console \code{>} belongs to the shell: it sends a command's output to a file. \code{select cone
150>60 10} is read as \code{select cone 150 10} with its output sent to \file{60}. In \file{/project}
the shell can't make that file and says \ui{geas: redirect: nothing can be made here: /60}; in
\file{/home} it makes it. Either way the command saw the wrong words. Type the arrow (\code{150→60}) or
quote the attitude (\code{'150>60'}), and \code{rm 60} the stray file.

\subsection{\code{rm} refuses to delete data}

\ui{rm: data/tocher.csv: “tocher” is drawn on Net 1 — remove those layers first}, naming every
diagram that draws it. A dataset that a
diagram draws can't be removed from the console: remove its layers (\code{rm "diagrams/Net 1/layers/tocher.csv"})
and then the data. The tree's \ui{Delete dataset} asks you instead, and removes both.

\subsection{Ctrl+Z does nothing in the console}

While you type in the console, keys belong to the shell. Type \code{undo} (and \code{redo}), or click
outside the console and press \key{Ctrl}+\key{Z} (\cref{sec:undo}).

\subsection{The arrow keys don't move along the line}

The console's line editor appends and deletes at the end only: \key{←} and \key{→} do nothing, and
\key{↑} \key{↓} recall earlier lines. Press \key{Backspace} to the mistake, or recall the line and retype
its end. \key{Ctrl}+\key{C} abandons a line.

\section{Projects and files}\label{sec:trouble-projects}

\subsection{My last session didn't come back}

The desk keeps a copy of the project between saves (\cref{sec:autosave}), and when it restores one the
status line says \ui{restored your last session (…)}. It didn't, because:

\begin{itemize}
\item the window was private, or the browser's data for the site was cleared;
\item wuffle opened from a different place: the file on your desktop and the website keep separate
copies (\cref{sec:where-kept});
\item it is a page that carries a project: each such page keeps its own copy, apart from the plain
wuffle's;
\item \menu{File>New project} was chosen, which forgets the kept copy.
\end{itemize}

The kept copy is a convenience. Save the project for anything you mean to keep.

\subsection{Save as page is greyed out}

It works from the published \file{wuffle.html}, which carries the page it copies. A wuffle run from its
source folder, as developers do, has no single page to copy (\cref{sec:save-page}).

\subsection{A project opened \enquote{but} some things were skipped}

\ui{… opened, but:} lists each part wuffle couldn't use, such as \ui{diagrams/Net 1/layers/x.json:
its data (data/x.csv) isn't in the file}: a layer whose dataset was removed from the zip by hand. The rest
of the project is open. Put the missing file back in the zip, or draw the data again.

\subsection{The signature does not match}

wuffle opens the log anyway, keeps it exactly as it came, and says:

\begin{lstlisting}
wuffle-log.csv: the signature does not match.

the bytes are NOT what the key signed — the log was changed after signing,
or the .sig belongs to another log
\end{lstlisting}

Either the CSV was edited (by a spreadsheet that \enquote{helpfully} resaved it, say), or the
\file{.sig} came from another publish. Find the copy the phone wrote, untouched, in its
\file{Downloads}, and open that. Don't treat a log that fails as the phone's record
(\cref{sec:field-logs}).

\subsection{The log says \enquote{unsigned}}

No \file{.sig} came with it, as far as wuffle could tell. The signature must be opened with its log and
named after it: \file{wuffle-log.csv.sig} beside \file{wuffle-log.csv}. A browser that renames a second
download (\file{wuffle-log (1).csv}) may rename one file and not the other; rename the \file{.sig} to
match its log, plus \file{.sig}. A \file{.sig} opened alone is ignored, and the status line says to open
it with its log. A log published outside the shell has no signature at all (\cref{sec:publish}).
